\section{Experiments}\label{sec:experiment}

\subsection{Baselines}\label{sec:baselines}

In order to study the dataset's characteristics, we implemented five baselines with different capabilities.
We briefly explain the five baselines that we implemented below, but more details can be found in \cref{sec:appendix-experiments-baseline}

\begin{description}
    \item[Majority vote] A baseline that outputs an oracle majority label for each hypothesis (NLI only).
    \item[Doc TF-IDF+SVM] A document-level multi-class linear Support Vector Machine \cite[SVM;][]{chang_libsvm_2011} with unigram bag-of-words features (NLI only).
    \item[Span TF-IDF+Cosine] Evidence identification based on unigram TF-IDF cosine similarities between each hypothesis and each span  (evidence identification only).
    \item[Span TF-IDF+SVM] A span-level binary Linear SVM with unigram bag-of-words features  (evidence identification only).
    \item[SQuAD BERT] A Transformer-based model as in the previous works discussed in \cref{sec:system-model}.
    Instead of allowing it to predict spans at arbitrary boundaries, we calculate a score for each of predefined spans by averaging token scores associated with the start and end of the span over different context windows.
    This makes sure that its performance is not discounted for getting span boundaries wrong.
\end{description}

\subsection{Experiment Settings}\label{sec:experiment-settings}

For evidence identification, we report mean average precision (mAP) that is micro averaged over labels.
We also report precision at recall 0.8 (P@R80) that is micro averaged over documents and labels.
P@R80 is the precision score when the threshold for evidence identification is adjusted to achieve a recall score of 0.8.
It was used in \cite{hendrycks_cuad_2021} to measure efficacy of a system under a required coverage level that is similar to typical human's.

\begin{table*}[t]
    \centering
    \begin{threeparttable}
        \centering
        \fontsize{8pt}{10pt}\selectfont
        \setlength{\tabcolsep}{4pt}
        \renewcommand{\arraystretch}{.6}
        \begin{tabular}{llc@{\hskip0pt}cc@{\hskip0pt}cc@{\hskip0pt}cc@{\hskip0pt}cc@{\hskip0pt}c}\toprule
            &  & \multicolumn{4}{c}{Evidence} & \multicolumn{6}{c}{NLI} \\\cmidrule(lr){3-6}\cmidrule(lr){7-12}
            Backbone Model & Fine-tuning Method & \multicolumn{2}{c}{mAP} & \multicolumn{2}{c}{P@R80} & \multicolumn{2}{c}{Acc.} & \multicolumn{2}{c}{F1 (C)} & \multicolumn{2}{c}{F1 (E)} \\\midrule
            BERT\textsubscript{\textit{base}} & None & .885 & \textsubscript{.025} & .663 & \textsubscript{.093} & .838 & \textsubscript{.020} & .287 & \textsubscript{.022} & .765 & \textsubscript{.035} \\
            BERT\textsubscript{\textit{large}} & None & .922 & \textsubscript{.006} & .793 & \textsubscript{.018} & .875 & \textsubscript{.006} & .357 & \textsubscript{.039} & .834 & \textsubscript{.002} \\
            DeBERTa v2\textsubscript{\textit{xlarge}} & None & .933 & \textsubscript{.002} & .859 & \textsubscript{.008} & .885 & \textsubscript{.001} & .360 & \textsubscript{.027} & .855 & \textsubscript{.002} \\
            \midrule
            BERT\textsubscript{\textit{base}} & Pretrained from scratch using a case law corpus \cite{zheng_when_2021} & .870 & \textsubscript{.015} & .578 & \textsubscript{.052} & .831 & \textsubscript{.032} & .289 & \textsubscript{.026} & .783 & \textsubscript{.040} \\
            BERT\textsubscript{\textit{base}} & Fine-tuned on case law and contract corpora \cite{chalkidis_legal-bert_2020} & .925 & \textsubscript{.004} & .811 & \textsubscript{.002} & .794 & \textsubscript{.008} & .272 & \textsubscript{.008} & .746 & \textsubscript{.018} \\
            DeBERTa v2\textsubscript{\textit{xlarge}} & Fine-tuned on span identification \cite{hendrycks_cuad_2021} & .936 & \textsubscript{.002} & .860 & \textsubscript{.003} & .892 & \textsubscript{.001} & .405 & \textsubscript{.016} & .859 & \textsubscript{.005} \\
            \midrule
            BERT\textsubscript{\textit{base}} & Fine-tuned on NDAs & .892 & \textsubscript{.002} & .690 & \textsubscript{.014} & .864 & \textsubscript{.004} & .326 & \textsubscript{.014} & .820 & \textsubscript{.010} \\
            BERT\textsubscript{\textit{large}} & Fine-tuned on NDAs & .922 & \textsubscript{.003} & .837 & \textsubscript{.008} & .875 & \textsubscript{.000} & .389 & \textsubscript{.009} & .839 & \textsubscript{.003} \\
            \bottomrule
        \end{tabular}
        \makeatletter\def\TPT@hsize{}\makeatletter
        \begin{tablenotes}[para,flushleft]
            \raggedright
            \fontsize{7pt}{7pt}\selectfont
            Refer to \cref{sec:experiment-settings} for the details on the metrics.
        \end{tablenotes}
    \end{threeparttable}
    \caption{Results for different backbone and pretrained models}\label{tab:results-pretrained}
\end{table*}

\begin{table}[t]
    \centering
    \begin{threeparttable}
        \centering
        \fontsize{8pt}{10pt}\selectfont
        \setlength{\tabcolsep}{2pt}
        \renewcommand{\arraystretch}{.6}
        \begin{tabular}{lc@{\hskip0pt}cc@{\hskip0pt}cc@{\hskip0pt}cc@{\hskip0pt}cc@{\hskip0pt}c}\toprule
            & \multicolumn{4}{c}{Evidence} & \multicolumn{6}{c}{NLI} \\\cmidrule(lr){2-5}\cmidrule(lr){6-11}
            & \multicolumn{2}{c}{mAP} & \multicolumn{2}{c}{P@R80} & \multicolumn{2}{c}{Acc.} & \multicolumn{2}{c}{F1 (C)} & \multicolumn{2}{c}{F1 (E)} \\\midrule
            Majority vote & \multicolumn{2}{c}{---} & \multicolumn{2}{c}{---} & .674 &  & .083 &  & .428 & \\
            Doc TF-IDF+SVM & \multicolumn{2}{c}{---} & \multicolumn{2}{c}{---} & .733 &  & .197 &  & .641 & \\
            Random & .024 & & .000 &  & \multicolumn{2}{c}{---} & \multicolumn{2}{c}{---} & \multicolumn{2}{c}{---} \\
            Span TF-IDF+Cosine & .381 & & .057  &  & \multicolumn{2}{c}{---} & \multicolumn{2}{c}{---} & \multicolumn{2}{c}{---} \\
            Span TF-IDF+SVM & .836 & & .322 &  & \multicolumn{2}{c}{---} & \multicolumn{2}{c}{---} & \multicolumn{2}{c}{---} \\
            SQuAD (BERT\textsubscript{\textit{base}}) & .825 & \textsubscript{.004} & .574 & \textsubscript{.004} & \multicolumn{2}{c}{---} & \multicolumn{2}{c}{---} & \multicolumn{2}{c}{---} \\
            SQuAD (BERT\textsubscript{\textit{large}}) & .869 & \textsubscript{.005} & .661 & \textsubscript{.043} & \multicolumn{2}{c}{---} & \multicolumn{2}{c}{---} & \multicolumn{2}{c}{---} \\
            \midrule
            Ours (BERT\textsubscript{\textit{base}})  & .885 & \textsubscript{.025} & .663 & \textsubscript{.093} & .838 & \textsubscript{.020} & .287 & \textsubscript{.022} & .765 & \textsubscript{.035} \\
            Ours (BERT\textsubscript{\textit{large}}) & .922 & \textsubscript{.006} & .793 & \textsubscript{.018} & .875 & \textsubscript{.006} & .357 & \textsubscript{.039} & .834 & \textsubscript{.002} \\
            \bottomrule
        \end{tabular}
        \makeatletter\def\TPT@hsize{}\makeatletter
        \begin{tablenotes}[para,flushleft]
            \raggedright
            \fontsize{7pt}{7pt}\selectfont
            Refer to \cref{sec:experiment-settings} for the details on the metrics.
        \end{tablenotes}
    \end{threeparttable}
    \caption{Main results}\label{tab:results}
\end{table}


For NLI, we report accuracy, a F1 score for contradiction (F1 (C)) and for entailment (F1 (E)).
We micro average these scores over documents and then macro average over labels.
This is to avoid the label imbalance to \emph{cancel out} with micro averaging and the results to appear too optimistic.

For our Span NLI BERT, we ran the same experiment ten times with different hyperparameters (detailed in \cref{sec:appendix-experiments-hyperparameters}) and report the average score of three models with the best development scores.
Since NLI is more challenging than evidence identification, we used macro average NLI accuracy for the criterion.
For the SQuAD BERT baseline, we ran hyperparameter search over 18 hyperparameter sets as described in \cite{devlin_bert_2019} and likewise report the average score of the three best models.
The metrics for the experiments with the hyperparameter search are followed by subscript numbers each of which denotes standard deviation of metrics over three runs.

\subsection{Results}

We first compared Span NLI BERT against baselines (\cref{tab:results}).
Span NLI BERT performed significantly better than the baselines, both in terms of evidence identification and NLI.
Nevertheless, the performance for contradiction labels is much worse than that of entailment labels, due to the imbalanced label distribution.
In terms of evidence identification, SQuAD BERT's mAP score was no better than that of Span TF-IDF+SVM, which illustrates the importance of explicitly incorporating span boundaries to input.

We then compared Span NLI BERT's performance with different backbone models and pretraining corpora including DeBERTa v2 \cite{he_deberta_2021} which was most successful in \cite{hendrycks_cuad_2021} (\cref{tab:results-pretrained}).
We can observe that making the models bigger benefits both evidence identification and NLI.
Fine-tuning models on legal corpora had mixed results.
Using a model pretrained on a case law corpus \cite{zheng_when_2021} did not benefit evidence identification nor NLI.
Fine-tuning BERT\textsubscript{\textit{base}} on NDAs has slightly improved the performance but the benefit is no longer visible for BERT\textsubscript{\textit{large}}.
Transferring DeBERTa\textsubscript{\textit{xlarge}} trained on CUAD \cite{hendrycks_cuad_2021} gave marginal improvement on NLI, making it the best performing model on the ContractNLI dataset.

\begin{table}[t]
    \centering
    \begin{threeparttable}
        \centering
        \fontsize{8pt}{10pt}\selectfont
        \setlength{\tabcolsep}{2pt}
        \renewcommand{\arraystretch}{.6}
        \begin{tabular}{lc@{\hskip0pt}cc@{\hskip0pt}cc@{\hskip0pt}cc@{\hskip0pt}cc@{\hskip0pt}c}\toprule
            & \multicolumn{4}{c}{Evidence} & \multicolumn{6}{c}{NLI} \\\cmidrule(lr){2-5}\cmidrule(lr){6-11}
            Hypothesis Usage & \multicolumn{2}{c}{mAP} & \multicolumn{2}{c}{P@R80} & \multicolumn{2}{c}{Acc.} & \multicolumn{2}{c}{F1 (C)} & \multicolumn{2}{c}{F1 (E)} \\\midrule
            Symbol (BERT\textsubscript{\textit{base}})  & .857 & \textsubscript{.044} & .574 & \textsubscript{.136} & .830 & \textsubscript{.014} & .294 & \textsubscript{.075} & .751 & \textsubscript{.027} \\
            Symbol (BERT\textsubscript{\textit{large}}) & .894 & \textsubscript{.020} & .703 & \textsubscript{.092} & .849 & \textsubscript{.006} & .303 & \textsubscript{.058} & .794 & \textsubscript{.026} \\
            \midrule
            Text (BERT\textsubscript{\textit{base}})  & .885 & \textsubscript{.025} & .663 & \textsubscript{.093} & .838 & \textsubscript{.020} & .287 & \textsubscript{.022} & .765 & \textsubscript{.035} \\
            Text (BERT\textsubscript{\textit{large}}) & .922 & \textsubscript{.006} & .793 & \textsubscript{.018} & .875 & \textsubscript{.016} & .357 & \textsubscript{.039} & .834 & \textsubscript{.002} \\
            \bottomrule
        \end{tabular}
        \makeatletter\def\TPT@hsize{}\makeatletter
        \begin{tablenotes}[para,flushleft]
            \raggedright
            \fontsize{7pt}{7pt}\selectfont
            Refer to \cref{sec:experiment-settings} for the details on the metrics.
        \end{tablenotes}
    \end{threeparttable}
    \caption{A controlled experiment using a randomly initialized special token for each hypothesis (Symbol) instead of hypothesis' surface tokens (Text)}\label{tab:results-symbolic}
\end{table}

